\documentclass[11pt]{article}
\usepackage{amsmath,amssymb,amsthm}
\usepackage{geometry}
\geometry{margin=1in}
\newtheorem{lemma}{Lemma}
\newtheorem{theorem}{Theorem}

\title{Havelock Lagrangian Eigenvalues:\\ First-Principles Derivation}
\author{Sandbox for Paper~I, Item~\#1}
\date{\today}

\begin{document}
\maketitle

\section{Setup}\label{sec:setup}

Let $z_k = e^{i\theta_k}$, $\theta_k = 2\pi k / N$, $k = 0, \ldots, N-1$,
be the regular $N$-gon of unit-circulation point vortices on the unit
circle. The Thomson energy is
\begin{equation}\label{eq:thomson}
  H = -\sum_{j<k} \ln |z_j - z_k|.
\end{equation}
We perturb each vortex in its \emph{local} (radial, tangential) frame:
\begin{equation}\label{eq:perturb}
  z_k \;\longmapsto\; z_k\bigl(1 + \delta r_k\bigr) + i\,z_k\,\delta t_k,
\end{equation}
with $\delta r_k, \delta t_k \in \mathbb{R}$ small. The $\mathbb{Z}_N$
cyclic symmetry $k \mapsto k+1 \pmod N$ permutes the vortices and
commutes with~$\nabla^2 H$, so the Hessian block-diagonalizes in the
discrete Fourier basis. The mode-$m$ cosine sector is parametrized by
\begin{equation}\label{eq:mode-m}
  \delta r_k = a\,\cos(2\pi m k / N), \qquad
  \delta t_k = b\,\cos(2\pi m k / N),
  \quad m \in \{1, \ldots, N-1\},
\end{equation}
(modes $m$ and $N-m$ yield identical blocks — the palindromic pair).
The corresponding $2 \times 2$ Hessian block in the $(a, b)$ basis is
what we compute.

\section{Two auxiliary identities}\label{sec:aux}

Two closed-form trigonometric sums enter the block computation and are
recorded here. Both are verified to 30 decimal digits in
\texttt{numerical\_check.py} (\texttt{check\_havelock\_identity\_rr},
\texttt{check\_cosecant\_squared\_sum}).

\begin{lemma}[Havelock identity]\label{lem:havelock}
For all integers $N \ge 2$ and $m \in \{1, \ldots, N-1\}$,
\begin{equation}\label{eq:havelock}
  S(m) \;:=\; \sum_{p=1}^{N-1} \frac{1 - \cos(2\pi m p / N)}{4 \sin^2(\pi p / N)}
  \;=\; \frac{m(N-m)}{2}.
\end{equation}
\end{lemma}

\begin{proof}
Using $1 - \cos(2\pi m p / N) = 2 \sin^2(\pi m p / N)$,
\[
  S(m) = \sum_{p=1}^{N-1} \frac{\sin^2(\pi m p / N)}{2 \sin^2(\pi p / N)}.
\]
The Dirichlet-kernel identity
$\sum_{p=1}^{N-1} \sin^2(\pi m p / N) / \sin^2(\pi p / N) = m(N - m)$
for $1 \le m \le N - 1$ (provable via the Chebyshev expansion
$\sin(m x) / \sin(x) = U_{m-1}(\cos x)$ and orthogonality of the
$U$-polynomials; see e.g.\ Gradshteyn and Ryzhik, Formula 1.342)
yields the claim.
\end{proof}

\begin{lemma}[Cosecant-squared sum]\label{lem:csc}
For $N \ge 2$,
\begin{equation}\label{eq:csc}
  \sum_{p=1}^{N-1} \frac{1}{4 \sin^2(\pi p / N)}
     \;=\; \frac{N^2 - 1}{12}.
\end{equation}
\end{lemma}

\begin{proof}
Standard; follows from partial-fraction expansion of $\cot(N z)$ and
taking the limit $z \to 0$, or equivalently from differentiating
$\prod_{p=1}^{N-1} \sin(\pi p / N) = N / 2^{N-1}$.
\end{proof}

\begin{corollary}[Cosine-weighted inverse-sine sum]\label{cor:cos-over-sin}
For $N \ge 2$ and any integer $k$,
\begin{equation}\label{eq:cos-over-sin}
  C(k) \;:=\; \sum_{p=1}^{N-1} \frac{\cos(2\pi k p / N)}{4 \sin^2(\pi p / N)}
       \;=\; \frac{N^2 - 1}{12} - \frac{k(N-k)}{2}
       \quad \text{for } 0 \le k \le N.
\end{equation}
The expression is periodic in $k$ modulo $N$, i.e.\ $C(k + N) = C(k)$.
\end{corollary}

\begin{proof}
Subtract \eqref{eq:havelock} from \eqref{eq:csc}:
$C(k) = (N^2-1)/12 - S(k) = (N^2-1)/12 - k(N-k)/2$.
Periodicity is immediate from the integrand.
\end{proof}

\section{Radial-radial block}\label{sec:rr}

Under a purely radial mode-$m$ perturbation
$\delta r_k = a\,\alpha_k$ with $\alpha_k = \cos(2\pi m k / N)$
(and $b = 0$ in \eqref{eq:mode-m}), the squared pair distance expands as
\begin{equation}\label{eq:pair-expand}
  |z_j(a) - z_k(a)|^2
  = |z_j - z_k|^2 \bigl[1 + a(\alpha_j + \alpha_k)\bigr]
    + a^2\,Q_{jk} + O(a^3),
\end{equation}
with
\[
  Q_{jk} := \alpha_j^2 + \alpha_k^2
            - 2 \alpha_j \alpha_k \cos(\theta_j - \theta_k),
  \qquad
  |z_j - z_k|^2 = 4 \sin^2\!\bigl(\pi(j-k)/N\bigr).
\]
(To see this, write $z_k(a) = z_k(1 + a\alpha_k)$, compute
$|z_j(a) - z_k(a)|^2 = |z_j - z_k|^2 + 2 a \operatorname{Re}\bigl[\overline{(z_j-z_k)}(\alpha_j z_j - \alpha_k z_k)\bigr] + a^2 |\alpha_j z_j - \alpha_k z_k|^2$,
and simplify the real part using
$\operatorname{Re}[(e^{-i\theta_j} - e^{-i\theta_k})(\alpha_j e^{i\theta_j} - \alpha_k e^{i\theta_k})]
 = (\alpha_j + \alpha_k)(1 - \cos(\theta_j - \theta_k))$.)

Expanding $-\ln|z_j(a) - z_k(a)| = -\tfrac12 \ln(|z_j(a) - z_k(a)|^2)$ to
second order in $a$ via $\ln(1 + x) = x - x^2/2 + O(x^3)$,
\begin{equation}\label{eq:log-expand}
  -\ln|z_j(a) - z_k(a)|
   = -\ln|z_j - z_k|
     - \tfrac{a}{2}(\alpha_j + \alpha_k)
     + a^2\!\left[\tfrac{(\alpha_j + \alpha_k)^2}{4}
                  - \tfrac{Q_{jk}}{2 |z_j - z_k|^2}\right] + O(a^3).
\end{equation}
The linear term, summed over pairs $j < k$, gives
$-(a/2)(N-1)\sum_k \alpha_k = 0$ since $\sum_k \cos(2\pi m k/N) = 0$ for
$m \in \{1, \ldots, N-1\}$ (equilibrium).

Write $\beta_m$ for the coefficient of $a^2$ in $H(a) - H(0)$:
\[
  \beta_m = \sum_{j<k}\!\left[\frac{(\alpha_j + \alpha_k)^2}{4}
                              - \frac{Q_{jk}}{2 |z_j - z_k|^2}\right].
\]
The unit-normalized radial Fourier mode vector $\hat{e}_r$ (in the
$2N$-dimensional Cartesian basis) has squared norm
$\|\hat{e}_r\|^2 = \sum_k \alpha_k^2 = N/2$ (valid for
$1 \le m \le N-1$, $m \ne N/2$; the case $m = N/2$ is handled
symmetrically and gives the same final formula by continuity).
The Hessian entry in this normalized direction is
\begin{equation}\label{eq:rr-norm}
  (\nabla^2 H)_{rr}^{(m)}
     = \frac{2\,\beta_m}{\|\hat{e}_r\|^2}
     = \frac{4\,\beta_m}{N}
     = \frac{2}{N}\sum_{j \ne k}\!\left[\frac{(\alpha_j + \alpha_k)^2}{4}
                                        - \frac{Q_{jk}}{2 |z_j - z_k|^2}\right].
\end{equation}
(The factor $2$ comes from $H''(0) = 2 \cdot (\text{coefficient of }a^2)$;
the factor $2$ converting $\sum_{j<k}$ to $\sum_{j \ne k}$ cancels half of it.)

Parametrize pairs by $p = k - j \pmod N$ with $p \in \{1, \ldots, N-1\}$
and apply the $N$-fold cyclic symmetry. Two elementary sums are needed:
\[
  \sum_{j=0}^{N-1} \alpha_j^2 = \frac{N}{2},
  \qquad
  \sum_{j=0}^{N-1} \alpha_j \alpha_{j+p} = \frac{N}{2}\cos(2\pi m p / N).
\]
From these,
\[
  \sum_{j=0}^{N-1} (\alpha_j + \alpha_{j+p})^2
     = N \bigl(1 + \cos(2\pi m p / N)\bigr),
  \qquad
  \sum_{j=0}^{N-1} Q_{j, j+p}
     = N \bigl(1 - \cos(2\pi m p / N) \cos(2\pi p / N)\bigr).
\]
Substituting into \eqref{eq:rr-norm},
\begin{equation}\label{eq:rr-intermediate}
  (\nabla^2 H)_{rr}^{(m)}
     = \sum_{p=1}^{N-1} \frac{1 + \cos(2\pi m p / N)}{2}
       \;-\; \sum_{p=1}^{N-1} \frac{1 - \cos(2\pi m p / N)\cos(2\pi p / N)}
                                   {4 \sin^2(\pi p / N)}.
\end{equation}

The first sum is elementary:
$\sum_{p=1}^{N-1} (1 + \cos(2\pi m p / N))/2 = (N-1)/2 - 1/2 = (N-2)/2$,
since $\sum_{p=1}^{N-1} \cos(2\pi m p / N) = -1$ for $m \in \{1, \ldots, N-1\}$.

The second sum reduces via the product-to-sum identity
$\cos A \cos B = \tfrac12[\cos(A-B) + \cos(A+B)]$:
\[
  1 - \cos(2\pi m p / N)\cos(2\pi p / N)
     = 1 - \tfrac12 \cos(2\pi(m-1) p / N) - \tfrac12 \cos(2\pi(m+1) p / N).
\]
Applying Lemma~\ref{lem:csc} and Corollary~\ref{cor:cos-over-sin},
\begin{align*}
  \sum_{p=1}^{N-1}\!\frac{1 - \cos(2\pi m p / N)\cos(2\pi p / N)}{4 \sin^2(\pi p / N)}
  &= \frac{N^2 - 1}{12}
     - \frac{1}{2}\Bigl[\frac{N^2 - 1}{12} - \frac{(m-1)(N-m+1)}{2}\Bigr] \\
  &\phantom{=} - \frac{1}{2}\Bigl[\frac{N^2 - 1}{12} - \frac{(m+1)(N-m-1)}{2}\Bigr] \\
  &= \frac{(m-1)(N-m+1) + (m+1)(N-m-1)}{4}
   = \frac{m(N-m) - 1}{2}.
\end{align*}
Combining,
\begin{equation}\label{eq:rr-unshifted}
  (\nabla^2 H)_{rr}^{(m)}
     = \frac{N - 2}{2} - \frac{m(N-m) - 1}{2}
     = \frac{N - 1 - m(N-m)}{2}.
\end{equation}
This is the radial-radial block eigenvalue of the \emph{unshifted}
energy Hessian $\nabla^2 H$. The Lagrange shift is applied in
\S\ref{sec:lag-shift}.

\end{document}
